\section{Tabel Parsing LL(1) dan Predictive Parser}
\label{sec:ch10-ll1}

\subsection{Konstruksi Tabel}

Tabel parsing \(M[A,a]\) berisi aturan produksi untuk nonterminal
\(A\) dan token lookahead \(a\). Konstruksinya:
\begin{enumerate}[label=\arabic*.,leftmargin=*]
  \item untuk setiap aturan \(A \to \alpha\) dan setiap
        \(a \in \text{FIRST}(\alpha)\), isi
        \(M[A,a] = A \to \alpha\);
  \item jika \(\varepsilon \in \text{FIRST}(\alpha)\), untuk setiap
        \(b \in \text{FOLLOW}(A)\), isi
        \(M[A,b] = A \to \alpha\);
  \item sel yang tidak terisi berarti kesalahan sintaksis.
\end{enumerate}

\begin{table}[H]
\centering
\caption{Tabel LL(1) untuk grammar ekspresi}
\label{tab:ch10-ll1-table}
\begin{tabular}{@{}l l l l l@{}}
\toprule
 & \code{id} & \code{+} & \code{*} & \code{\$} \\
\midrule
\(E\)  & \(E \to TE'\) & & & \\
\(E'\) & & \(E' \to +TE'\) & & \(E' \to \varepsilon\) \\
\(T\)  & \(T \to FT'\) & & & \\
\(T'\) & & \(T' \to \varepsilon\) & \(T' \to *FT'\) &
\(T' \to \varepsilon\) \\
\(F\)  & \(F \to \code{id}\) & & & \\
\bottomrule
\end{tabular}
\end{table}

Tabel pada \tabref{tab:ch10-ll1-table} belum memuat kolom
\code{(} dan \code{)}. Untuk kelengkapan, sel \(E\) dan \(T\)
pada kolom \code{(} diisi aturan yang sama, sedangkan sel
\(E'\) dan \(T'\) pada kolom \code{)} diisi \(\varepsilon\).

\subsection{Algoritma Predictive Parser}

Predictive parser menggunakan stack eksplisit:
\begin{enumerate}[label=\arabic*.,leftmargin=*]
  \item dorong \(\$\) dan simbol awal ke stack;
  \item \(a\) adalah token saat ini;
  \item selama stack tidak kosong:
        \begin{enumerate}[label=(\alph*),leftmargin=*]
          \item \(X\) adalah puncak stack;
          \item jika \(X\) terminal atau \(\$\), cocokkan dengan
                \(a\); jika cocok, pop dan majukan \(a\);
          \item jika \(X\) nonterminal, lihat \(M[X,a]\);
          \item jika aturan \(X \to Y_1\cdots Y_k\) ada, pop \(X\)
                dan dorong \(Y_k,\ldots,Y_1\);
          \item jika sel kosong, laporkan kesalahan.
        \end{enumerate}
\end{enumerate}

\subsection{Contoh Penelusuran}

Untuk masukan \(\code{id} + \code{id} * \code{id}\ \$\),
penelusuran stack adalah:
\begin{table}[H]
\centering
\caption{Penelusuran predictive parser}
\label{tab:ch10-trace}
\begin{tabularx}{\textwidth}{@{}l X l@{}}
\toprule
\textbf{Stack} & \textbf{Masukan} & \textbf{Aksi} \\
\midrule
\(\$E\) & \(\code{id} + \code{id} * \code{id}\ \$\) &
\(E \to TE'\) \\
\(\$E'T\) & \(\code{id} + \code{id} * \code{id}\ \$\) &
\(T \to FT'\) \\
\(\$E'T'F\) & \(\code{id} + \code{id} * \code{id}\ \$\) &
\(F \to \code{id}\) \\
\(\$E'T'\code{id}\) & \(\code{id} + \code{id} * \code{id}\ \$\) &
cocokkan \code{id} \\
\(\$E'T'\) & \(+ \code{id} * \code{id}\ \$\) &
\(T' \to \varepsilon\) \\
\(\$E'\) & \(+ \code{id} * \code{id}\ \$\) &
\(E' \to +TE'\) \\
\(\$E'T+\) & \(+ \code{id} * \code{id}\ \$\) & cocokkan \(+\) \\
\(\$E'T\) & \(\code{id} * \code{id}\ \$\) & \(T \to FT'\) \\
\(\$E'T'F\) & \(\code{id} * \code{id}\ \$\) & \(F \to \code{id}\) \\
\(\$E'T'\code{id}\) & \(\code{id} * \code{id}\ \$\) &
cocokkan \code{id} \\
\(\$E'T'\) & \(* \code{id}\ \$\) & \(T' \to *FT'\) \\
\(\$E'T'F*\) & \(* \code{id}\ \$\) & cocokkan \(*\) \\
\(\$E'T'F\) & \(\code{id}\ \$\) & \(F \to \code{id}\) \\
\(\$E'T'\code{id}\) & \(\code{id}\ \$\) & cocokkan \code{id} \\
\(\$E'T'\) & \(\$\) & \(T' \to \varepsilon\) \\
\(\$E'\) & \(\$\) & \(E' \to \varepsilon\) \\
\(\$\) & \(\$\) & terima \\
\bottomrule
\end{tabularx}
\end{table}

Penelusuran pada \tabref{tab:ch10-trace} menunjukkan bahwa parser
tidak pernah mundur. Setiap keputusan diambil dari tabel berdasarkan
token saat ini, sehingga waktu parsing linear terhadap panjang
masukan.
